\chapter{Down Syndrome Tests}
\section*{What Is This Test?} Down syndrome testing includes prenatal screening, prenatal diagnostic testing, and postnatal chromosome analysis for trisomy 21.
\section*{Alternative Names} Trisomy 21 testing; Down syndrome screening; chromosome 21 testing.
\section*{Principle} Screening estimates risk using ultrasound and maternal blood or cell-free DNA; diagnostic testing examines fetal or newborn cells.
\section*{Why This Test Is Ordered} It is used during pregnancy for risk assessment or diagnosis and after birth when physical findings or screening suggest Down syndrome.
\section*{Primary vs Secondary Test} Screening is preliminary. CVS or amniocentesis is diagnostic prenatally; newborn karyotype or chromosomal microarray confirms postnatally.
\section*{How This Test Is Performed} Samples may be maternal blood, ultrasound measurements, placental tissue, amniotic fluid, or newborn blood.
\section*{What Preparation Is Needed} Usually none; genetic counseling is recommended before diagnostic testing.
\section*{Understanding Results} A positive screen is not a diagnosis. A diagnostic result must be interpreted with counseling and clinical findings.
\section*{Follow-up Tests} Detailed ultrasound, fetal echocardiography, hearing assessment, thyroid testing, blood count, and developmental evaluation may be recommended.
\section*{References} \begin{enumerate}\item MedlinePlus. Down Syndrome Tests.\item American College of Obstetricians and Gynecologists. Prenatal genetic testing guidance.\end{enumerate}
\clearpage\section*{Understanding Results and Follow-up}
The laboratory report should identify the specimen, method, units, reference interval or decision limit, and whether the result is positive, negative, abnormal, or inconclusive. Reference intervals vary with age, sex, pregnancy, specimen type, assay, and laboratory; a result outside a range is not automatically a diagnosis, and a result within a range does not always exclude disease. Discuss unexpected results with the ordering clinician, who can decide whether repeat or confirmatory testing, treatment, or referral is appropriate. Do not change medicines or delay urgent care based on an isolated result.
\section*{Regulatory Status and Authorization Date}This chapter describes a test category, not a specific commercial product. FDA, CDSCO, EU IVDR, and MHRA approval or registration dates therefore cannot be assigned without the assay manufacturer, product name, instrument, and intended use. For laboratory-developed tests, an individual agency approval date may not exist. See the Preface for the interpretation of regulatory status.
\clearpage
